%==============================================================================
% Homework Assignment Template
%==============================================================================
\documentclass[11pt]{article}

%------------------------------------------------------------------------------
% PACKAGES
%------------------------------------------------------------------------------
\usepackage[margin=1in]{geometry}   % Adjust margins
\usepackage{amsmath,amssymb,amsthm} % Math packages
\usepackage{graphicx}               % Include images
\usepackage{hyperref}               % Hyperlinks
\usepackage{enumitem}               % Better control over lists
\usepackage{fancyhdr}               % Header/footer customization

%------------------------------------------------------------------------------
% CUSTOM COMMANDS & SETTINGS
%------------------------------------------------------------------------------
\newcommand{\courseName}{Ensemble Quantum Computing Laborator}     % Change to your course name
\newcommand{\assignmentTitle}{Prelab 0} % Change to your assignment title
\newcommand{\studentName}{Runzhao Guo}       % Change to your name
\newcommand{\Date}{January 12, 2025}   % Change to your due date

% Set up the page header
\pagestyle{fancy}
\fancyhf{} % Clear header and footer
\lhead{\courseName}
\chead{\assignmentTitle}
\rhead{\studentName}
\rfoot{\thepage}



%------------------------------------------------------------------------------
% DOCUMENT START
%------------------------------------------------------------------------------
\begin{document}

%------------------------------------------------------------------------------
% TITLE SECTION
%------------------------------------------------------------------------------
\begin{center}
    {\large \textbf{\courseName}} \\
    \vspace{3pt}
    {\Large \textbf{\assignmentTitle}} \\
    \vspace{3pt}
    \studentName \\
    \vspace{3pt}
    \Date
\end{center}

\vspace{0.5in} % Add some space before the content starts

%------------------------------------------------------------------------------
% INSTRUCTIONS / ABSTRACT
%------------------------------------------------------------------------------
%\noindent \textit{Instructions:} Provide any useful instructions or an abstract 
% for your homework assignment here. For example, due date, collaboration policies, 
% citation requirements, etc.

%------------------------------------------------------------------------------
% PROBLEM 1
%------------------------------------------------------------------------------
\section*{Values and dimensions of relevance to NMR}

\section{} 
\begin{enumerate}[label=\alph*)]
    \item Magnetic moment of a Dirac electron: 
    \[
        \mu_{Dirac} = \frac{e \hbar}{2m} = 9.274010 \times 10^{-27} \text{J/T}.
    \]
    \[
        1\,\text{eV} = 1.602176634 \times 10^{-19}\,\text{J}.
    \]
    \[
        \mu_{\mathrm{Dirac}} = 5.788 \times 10^{-8}\,\text{eV}.
    \]


    \item The current (CODATA-recommended) best-measured value for the magnetic moment of an electron is $9.2848 \times 10^{-27} \text{J/T} = 5.79 \times10^{-8} \text{eV/T}$
    \item A simple expectation of the magnetic moment of a proton using the same method is:
    \[
        \mu_{Proton} = \frac{e \hbar}{2m_{Proton}} = 5.0508 \times 10^{-27} \text{J/T} = 3.152 \times 10 ^{-8} \text{eV/T}
    \]
    \item However the measured value is very different:
    \[
        \mu_{p measured} = 1.4106 \times 10^{-26} \text{J/T} = 8.8030 \times 10^{-8} \text{eV/T}
    \]
    Which is 2.7 times larger the naive expectation, the reason to this is the g fector of a proton is approximately 5.585, which is way bigger then 2.
\end{enumerate}

\section{}
\begin{enumerate}[label=\alph*)]
    \item 
        Using equaton $\gamma = \frac{2\mu}{\hbar} $, The gyromagnetic ratio of the expected magnetic moment of Dirac electon is $1.7589 \times10^{8} \text{rad} \cdot \text{s}^{-1} \cdot \text{T}^{-1}$.\\
        \\
        The gyromagnetic ration of the measured magnetic moment of Dirac electon is $1.7609 \times10^{8} \text{rad} \cdot \text{s}^{-1} \cdot \text{T}^{-1}$.\\ 
        \\
        The gyromagnetic ratio of the proton(exptected) is $9.5944 \times10^{7} \text{rad} \cdot \text{s}^{-1} \cdot \text{T}^{-1}$.\\
        \\
        The gyromagnetic ratio of the proton(measured) is $2.6753 \times10^{8} \text{rad} \cdot \text{s}^{-1} \cdot \text{T}^{-1}$.\\
    \item In \text{MHz/T}:
        The gyromagnetic ratio of the expected magnetic moment of Dirac electon is $2.8007 \times10^{7} \text{MHz/T}$.\\
        \\
        The gyromagnetic ration of the measured magnetic moment of Dirac electon is $2.803 \times10^{7} \text{MHz/T}.$\\ 
        \\
        The gyromagnetic ratio of the proton(exptected) is $1.5277 \times10^{7} \text{MHz/T}.$\\
        \\
        The gyromagnetic ratio of the proton(measured) is $4.26 \times10^{7} \text{MHz/T}$.\\
\end{enumerate}

\section{}
    The simplest reason is that the proton is not a point-like Dirac particle but a composite object made of quarks bound together by gluons. Its internal quark-gluon dynamics gives it a magnetic moment that deviates significantly from the value predicted by the naive Dirac equation.

\section{}
\begin{enumerate}[label=\alph*)]
    \item in Figure 1
    \begin{figure}[ht]
        \centering
        \includegraphics[width=0.5\textwidth]{/Users/runzhaoguo/Documents/MQST-UCLA/411/lab0/Figure/Figure_1.png}
        \caption{The energy level diagram for the spin-up and spin-down states of the proton as a function of the magnetic field.}
        \label{fig:1}
    \end{figure}

    \item 
    \[
    \triangle E_Z = E_{\frac{1}{2}} - E_{-\frac{1}{2}} = \gamma \hbar B
    \]

    \item The dependency of $\triangle E_Z $ on magnetic field, in MHz/T:
    from 
    \[
        \triangle E_Z = \gamma \hbar B
    \]
    we know, the dependency of $\triangle E_Z$ on magnetic field is $2.823 \times 10^{-26} $\text{MHz/T}

    \item let polarization be $P = \frac{N_{\uparrow} - N_{\downarrow}}{N_{\uparrow} + N_{\downarrow}}$.\\
    Assume Boltzmann distribution, $P_i = \frac{1}{Z} exp(-\frac{E_i}{k_B T})$, $P = tanh\frac{\gamma \hbar B}{2 k_B T}$.
\end{enumerate}

\section{}
Volume of the tube:
\[
    v_{tube} = \pi \times (\frac{d}{2})^2 \times l = 0.393\text{mL}
\]
volume of $H_2O$:
\[
    v_{H_2O} = 0.05 \times v_{tube} = 0.01965\text{mL}
\]
\begin{enumerate}[label=\alph*)]
    \item 
    \[
        m_{H_2O} = 0.01965\text{g}
    \]
    \[
        n_{H_2O} = \frac{0.01965\text{g}}{18\text{g/mol}} = 1.09 \times 10^{-3} \text{mol}
    \]
    \[
        N_{H_2O} = n_{H_2O} \times N_A = 6.56 \times 10^{20} \text{molecules}
    \]
    which means $1.312 \times 10^{21}$ Hydrogen nuclei, that is $ 3.3 \times 10^{21}$ Hydrogen nuclei per mL.
    \item $N_{H} = 1.312 \times 10^{21}$
    \item $2.18 \times 10^{-3} \text{mol}$
\end{enumerate}
\section*{Larmor frequency, the rotating frame, and offset frequencies }
\section{}
\begin{enumerate}[label=\alph*)]
    \item the number of protons in spin up/down states can be calculated as such:
    \[
        \triangle N = P N_H = 1.367 \times 10^{16}
    \]
    \[
        N_{\uparrow} =\frac{1}{2} (1+P) N_H = 6.560 \times 10^{20}
    \]
    \[
        N_{\downarrow} =\frac{1}{2} (1-P) N_H = 6.559 \times 10^{20}
    \]

    \item the net magnetic moment can be calculated as such:
    \[
    \mu_{net} = \triangle N \mu_{proton} = 7.716 \times 10^{-11}
    \]

    \item the order of magnitude of the magnetic field strength at 2cm is $10^{-12}\text{T}$
\end{enumerate}

\section{}
the shifts are calculated as such:$\delta = 10^{6} \times \frac{\nu - \nu_{TMS}}{\nu_{TMS}}$, so in oreder to do get the physical value of the shifts, we do the revers of this calculation.
\[
    \triangle \delta = 3.80\text{ppm}
\]

\[
    \triangle \nu = \triangle \delta \times 400.13 = 1520.494 \text{Hz} \approx 9550 \text{rads/s}
\]

\section{}
The Larmor frequency is calculated as such: 
\[
    \nu = -\frac{1}{2\pi}\gamma(1+\delta)B
\]
using the conditions given in exercise, $\nu = -1.0072 \times 10^8 \text{Hz}$ or $ -6.3249 \times 10^{8}\text{rads/s}$

\section{}
What happened here is the RF pules rotates the net magnetization away from z-axis. The angle of rotation can be calculated as such:
\[
    \theta = \omega_1 t = 2\pi \nu_1 t
\]

From the data given, $\nu_1$ can be calculated as 
\[
    \nu_1 = \frac{\theta}{2 \pi t} \approx 24.4 \text{kHz}
\]
The $90$ degree pulse length is $10.25 \mu \text{s}$

As for the second null, it is clear that the magnetization has returned to +z.

\section{}
It does not. Here's why:
let's first get the $\omega_1$
    \[
        \omega_1 = \frac{\frac{\pi}{2}}{\tau_{pulse}} \approx 25 \text{kHz}
    \]
    The maximum offset here is:
    \[
        \triangle = 45 \text{MHz} \approx 2.83 \times 10^8 \text{rads/s}
    \]
    Notice $\triangle \gg \omega_1$, meaning the effective magnetic field is along the z-axis. The current magnitization can only tilte the spin by $\frac{\omega_1}{\triangle} \approx 1.4 \times 10^{-4}$ rads. Which is essentially negligible. Therefore the it does not have a significant effect on C-13 nuclei

\section{}
Before the added pules, the vector is along $+y$ direction. \\
$\phi = +x$, the vector will be pointing at $-z$ direction\\
$\phi = -x$, the vector will be pointgin at $+z$ direction\\
for both$\phi = -y$ and $\phi = +y$, the vector will be pointing at $+y$ direction. \\
% \begin{thebibliography}{9}
% \bibitem{exampleRef} Author Name, \textit{Title of the Book}, Publisher, Year.
% \end{thebibliography}

%------------------------------------------------------------------------------
% DOCUMENT END
%------------------------------------------------------------------------------
\end{document}
